\section{Introduction}
\label{sec:intro}
\parjanya{First para should be general introduction. Cite a few LLMs and how they have been integrated with tools. then say there are 1000s of tools and add citations supporting this.}

% P1: Context + problem. One idea per sentence.
Large language models are increasingly augmented with external tools to act as autonomous agents~\citep{schick2024toolformer, yao2023react, shen2024hugginggpt}. Most production agents today operate over curated toolsets where a human has already vetted each tool. But the ecosystem is shifting toward open marketplaces: standardized protocols like MCP~\citep{anthropic2025mcp} now expose thousands of community-contributed endpoints, and retrieval-augmented approaches already struggle at this scale---tool selection accuracy collapses from 43\% to under 14\% as registries grow~\citep{sun2025ragmcp}. The core difficulty is not finding tools that \emph{match} a query, but finding tools that \emph{work}. In open registries, multiple endpoints claim to serve the same function with near-identical descriptions, yet they differ sharply in reliability. We call this \emph{underdifferentiation}: tools that are semantically indistinguishable but qualitatively distinct.

\parjanya{Explain existing approaches properly - retrievers, trained reitrevers. explain downside.}

% P2: Why existing work fails. Argument first, citations as evidence.
No existing approach can resolve underdifferentiation, because all operate on static, description-level information. In curated settings this was sufficient---a human had already filtered for quality. In open marketplaces, quality variance is the default, and descriptions carry no signal about it. Retrieval methods rank by semantic similarity and are blind when descriptions converge~\citep{shi2024toolret}. Active discovery improves query formulation but still routes by similarity~\citep{fei2025mcpzero}. Recommendation approaches use static benchmark labels with no runtime signal~\citep{shi2026agentselect}. Deployed systems like Anthropic's tool search use BM25 over metadata~\citep{anthropic2025toolsearch}. The gap is consistent: none of these systems learn from whether a tool actually succeeded.

\parjanya{This should propose our method. You have already done this to some extent. But it should be clear how we solve previous issues: 1. Do not use fixed training data but learn online. Learn from LLM feedback. Explain what this feedback is in one line.  }
% P3: Our approach. High-level mechanism only; details go in Section 4.
We propose \textsc{Toolflix}, a recommender that learns which tools work from execution feedback. \textsc{Toolflix} places a learned reranker after a standard semantic retriever. The reranker has two heads: a \emph{relevance head} that learns query-tool affinity (e.g., fetch endpoints work for search tasks, but specific ones fail on certain domains), and a \emph{quality head} that learns per-tool reliability from accumulated success rates---the equivalent of a star rating that emerges from usage, not descriptions. A blending parameter $\alpha$ shifts from relevance-dominant at cold start to quality-dominant as feedback accumulates. After each task, an LLM judge provides a binary success/failure signal; training uses pairwise margin loss over these outcomes and requires no human labels.


\parjanya{More details on our experiments. Explain our experimental results in 1-2 lines. Can end with contribution.}
% P4: Key evidence preview. TODO: fill with final numbers.
% Testbed: 59 real + 104 synthetic MCP endpoints with controlled failure modes.
% Evaluation axes: (1) learning curve, (2) held-out comparison,
% (3) cold-start/OOD, (4) quality drift, (5) transfer to TOOLRET.
% Show underdifferentiation score correlates with reranker improvement.

% P5: Contributions. One sentence each.
Our contributions:
\begin{itemize}
    \item We formalize \emph{underdifferentiation}---high description similarity coupled with high variance in execution quality---and show that the reranker's gain over retrieval correlates with this measurable score.
    \item We propose a two-head reranker that separates learned relevance from learned quality, with an adaptive blending schedule for cold start.
    \item We release a synthetic tool benchmark with controlled quality variance: paraphrased descriptions and injected failure modes that are invisible to retrievers, isolating the value of feedback.
    \item We show feedback-driven reranking improves selection over retrieval baselines, with the largest gains in high-underdifferentiation categories.
\end{itemize}
\parjanya{These contributions make no sense. Our contribution is 1: showing existing retrivrs are inadequate 2. new recommendation system 3. our system can personalize to models 4. adapt to changing tools over time}